\section{Related Work}

Much of what this theory needs has already been proven by others, and the work mostly assembles established results onto one specific substrate, then asks one further question those programs leave aside, what the felt interior of the thing is and which way it moves. Honest credit, by family, is owed.

The idea of a discrete relational spacetime is old. The causal-set program of Bombelli, Lee, Meyer, and Sorkin, and of Rideout, Sorkin, and Surya, builds spacetime as a discrete order with distance as step-counting, buying Lorentz invariance with randomness where this work buys it with curvature. Causal dynamical triangulations, of Ambjorn, Jurkiewicz, and Loll, gets geometry and dimension as outputs of a discrete path integral, an ensemble average, where this work fixes one rigid geometry and never sums. Loop quantum gravity and spin networks make space a graph of relations, there a fluctuating quantum state, here a fixed graph whose only variable is the tone. The Wolfram Physics Project is the closest neighbour in spirit, a universe of graph plus local rule, but it searches a vast rule space and lets the graph rewrite itself, where this work fixes the graph by argument and varies only the tone.

That a deterministic classical substrate could underlie the quantum is 't Hooft's conviction, inherited rather than originated here. The route from a nearest-neighbour quantum walk to the Dirac equation, developed by D'Ariano and Perinotti, Bisio, Arrighi, and Farrelly and Short, is the construction this work leans on most heavily and does not claim, its contribution being the hyperbolic substrate, the forced geometry, and the demonstration that the knit is this construction. Holography and tensor networks, the work of Maldacena, Ryu and Takayanagi, Van Raamsdonk, Swingle's tensor networks, and the Pastawski-Yoshida-Harlow-Preskill code, are largely static, and the contribution here is to let a rule run and measure the area law and the working code as properties of the dynamics. The hyperbolic lattices of Kollar and of Maciejko and Rayan, and the hyperbolic cellular automata of Margenstern, supply the geometry and its universality, while Cannon and Hamann supply the geometric group theory that turns negative curvature into a speed limit and a holographic readout. The associative-computing paradigm of Potter is the model the memory engine realizes, made logarithmic here by the curvature.

On foundations and mind the debts are equally plain. The informational reconstructions of quantum theory by Chiribella, D'Ariano, and Perinotti derive its structure axiomatically, where this work commits to a substrate and lets it fall out. The hard problem of Chalmers, the integrated information of Tononi, and Wheeler's it-from-bit frame the experiential reading, which differs in taking experience as the substance rather than a quantity computed from structure. Entropic gravity, of Jacobson, Padmanabhan, and Verlinde, is the natural home for the result that $1/r$ gravity is emergent. Friston's active inference is, almost exactly, what a self here is, the forward model perception, the will action, the lean the preference. Hoel's causal emergence is the rigorous backbone for the degeneracy trick. Conway and Kochen's free-will theorems are the sharpest adversary to a deterministic base, named rather than dodged, and conceded at the micro level while freedom is claimed at the macro. The origin-of-life physics of England and of assembly theory, and the constructor theory of Deutsch and Marletto, frame the climb through life. Whitehead's process philosophy is the nearest classical ancestor of the whole.

What is genuinely new is narrow, and stated narrowly. A substrate argued to be selected rather than searched, rewritten, or assumed. The ternary tone with the lean as a built-in value direction, the single deliberate posit. One substrate aimed at physics and life and mind together. An experiential monism, the part furthest outside the literature. And a falsifiable single-substrate program that states its own negatives. The selection argument, the dimension ladder plus the spinor requirement, is the strongest of these, and the rest of the novelty leans on it.
